\section{边地首领如何成为西夏}

982年李继迁脱离宋朝羁縻控制，在西北边地集结力量。\eventref{XIA-0982-EVENT-049}{李继迁集结力量（982）}990年他在战与和之间接受宋廷封授，却维持独立军事行动。册封不是统治已经实现的同义词，也不是纯粹的外交辞令：它是围绕夏州、灵州、道路、部众与补给所作的暂时安排。宋方《夏国传》记录得最完整，却以敌手和朝贡的眼光组织材料；黑水城西夏文文书、河西考古和碑刻能补足制度与日常，却也有出土集中、年代与释读不均的限制。

1004年李继迁去世，李德明继承部众，1009年接受宋廷名号并经营夏州、灵州一带；1032年李元昊继位。继承并非单独的宫廷事件，它决定谁可以调用边地的军力、草场、城寨与贸易节点。1034年元昊改国号和年号、调整王权象征，1036年向甘、凉方向扩张，1038年称帝建国。宋朝的“僭”或“叛”是它的外交语言，不能代替党项统治者自身的政治工程，更不能用宋方战报推断西夏全境的族群认同或财政收入。

\section{宋夏战争与战后的政治}

1040年三川口、1041年好水川、1042年定川寨，宋军连续受挫；1044年双方订立庆历和议。宋史可确证战争与和议的骨架，却不能提供西夏军队规模和内部决策的完整自述。黄土高原的水源、山地道路、堡寨和补给，解释了何以一方难以迅速取胜；赐予、互市和称谓则是降低消耗的制度工具，不应被直接翻译为一方对另一方的永久从属。

1048年元昊遇害，1049年李谅祚即位，后族与权臣主政；1067年李秉常即位，幼主政治再现。后继者与摄政者的具体动机，在以宋方材料为主的叙事中并不透明，不能补写为确定的宫廷阴谋。可以看见的是，西夏必须在王室继承、边防与资源调配之间维持联盟；儿童君主并没有使这个国家暂停面对宋、辽和地方网络。

1080年宋神宗在西北进取，1082年永乐城被攻破。宋方叙述中的战败、伤亡与责任归属需要对读，西夏一侧的材料也不应由沉默而被想象为毫无代价。城寨是军事据点，也是取水、运粮、征发和驻军的持续负担；永乐城的结果因此属于长期边防财政，而非一位将领的单次得失。

\subsection{交读窗口：苏轼的“教战守”与永乐城}

传世本编次难以确证成篇年的《教战守策》，是北宋中期苏轼的政论散文，面向的是宋朝如何教民、守城和用兵的公共议论，而不是西北前线的战报。文中说：“然天下果未能去兵，则其一旦将以不教之民而驱之战，不败而已矣！”这句以设问和警语推进的文字，能使永乐城的取水、道路、守备与补给不再只是一场败仗的技术背景：在宋代士人的论说中，训练、地方组织和战争成本已经构成相连的问题。

不过，它讲的是一位宋人如何设想防务，不能说明西夏军队的训练、战略或一〇八二年攻城现场。永乐城的基本年序宜以《宋史》本纪、夏国传与《续资治通鉴长编》相校，《宋会要辑稿》的兵、蕃夷门可追问宋方制度文本；西夏方面则必须让黑水城西夏文文书、汉文与西夏文碑刻、城址和水利考古进入问题。后者多集中于保存较好的地点，文书纪年、释读和出土层位也不齐，不能用来补写战场细节。钱穆关于制度能否把中央意图转为边防能力的提问，在此须同吕思勉对征发、贸易和地方生计的追问并置；苏轼的策论可显示宋方焦虑，却不能让西夏成为只有宋方声音的对手。

\section{在金与蒙古之间的终局}

1086年李乾顺即位，1115年金建国改变北方外交环境，1123年西夏与金建立新的关系，1139年李仁孝即位，1144年双方以册封、贡赐和边境安排维持和平。此处的“册封”同样不能覆盖实际军事实力、关市交换和地方征收；应将名义、互惠与控制分别讨论。1193年李纯佑即位时，西夏仍在宋金之间经营自己的位置。

1206年蒙古诸部统一，1209年蒙古攻夏后出现和议与联姻安排，1217年双方军事要求已见紧张，1223年李德旺即位，1226年蒙古大举攻入，1227年西夏灭亡。蒙古、宋、金和后出的元代材料对战争的解释各不相同；城破、人口损失和接收范围尤其不能仅据征服者数字。西夏的终结改变了河西、宁夏和黑水城文书所见地区的统治者，却没有使语言、佛教、印刷、商路和居民经验一同消失。把这些材料作为主体证据，而不只把西夏写成宋或蒙古的边注，才是理解这个国家的必要条件。
